\item Suponga que un objeto tridimensional muy delgado tiene un espesor $dp$ a lo largo del eje óptico, de manera que se extiende desde la distancia objeto $p$ hasta $p + dp$.
\begin{enumerate}
    \item Demuestre que el espesor longitudinal de su imagen, $dq$, se conoce por $dq = -(q^2/p^2)dp$.
    \item Si el aumento longitudinal del objeto se define como $M_{\text{long}} = dq/dp$, ¿cómo se relaciona con la magnificación lateral $M$?
\end{enumerate}